\chapter{积分：把无数个小量加起来}

\section{里程表不量速度，却会算路程}

坐进一辆汽车，仪表盘上有两块表：一块是速度表，指针告诉你此刻开得多快；另一块是里程表，数字告诉你一共开了多远。上一章我们关心的是反方向的问题：知道里程表随时间的变化，怎样算出每一时刻的速度？答案是求导数——用很短一段时间里里程的增加量除以这段时间，再把时间段缩到无穷小。

现在把问题反过来。假设里程表坏了，只有速度表是好的，而且它每隔一秒就把当前速度记录下来。你能不能只用这些速度记录，算出这趟车一共跑了多远？

直觉上会这样做：第一秒里车速大约是每秒 $20$ 米，这一秒大约跑了 $20$ 米；第二秒车速大约是每秒 $22$ 米，这一秒大约跑了 $22$ 米……把每一秒跑的距离加起来，就是总路程的近似值。如果记录更勤快，每十分之一秒记一次，每段里的速度变化更小，加出来的结果就应该更准。

其实生活里到处都是这样一对一对的表：水龙头开多大是“率”，水表走了多少字是“量”；电器功率多大是“率”，电表走了多少度是“量”；网速多快是“率”，手机流量统计是“量”。每一块“总量表”都在默默做同一件事：把每时每刻的“率”连续不断地累加起来。本章就是要搞清楚，这种累加到底是什么运算。

这个朴素得不能再朴素的想法——\emph{把整个过程切成许多小段，在每小段里近似当成均匀的，再把所有小段的贡献加起来}——就是积分的全部起点。我们会看到，同一个想法能算出曲线围成的面积、球和圆锥的体积、变力做的功、随机事件的概率，还能回答一个更惊人的问题：无限多个数加起来，什么时候会得到一个有限的答案？

\section{一个把人卡住两千年的问题}

先从一个看起来更“几何”的问题开始。抛物线 $y=1-x^{2}$ 和 $x$ 轴围成一个像小山包一样的图形，$x$ 从 $-1$ 到 $1$。它的面积是多少？

三角形面积你会算，矩形面积你会算，圆的面积你背过公式。可是这个“山包”的边是弯的，既不是多边形也不是圆。凭直觉猜一猜：它外接的矩形宽 $2$、高 $1$，面积是 $2$；它肚子里能塞下的最大三角形底为 $2$、高为 $1$，面积是 $1$。所以答案在 $1$ 和 $2$ 之间。再猜得细一点？有人说 $1.5$，有人说接近 $1.3$，光靠直觉已经说不清了。

对付弯边图形，古人早有过辉煌的成功。圆的面积就是一个例子：魏晋时期的刘徽用“割圆术”算它——在圆内作内接正六边形，再作正十二边形、正二十四边形……边数每翻一倍，多边形就更贴近圆一圈。他的原话是：“割之弥细，所失弥少，割之又割，以至于不可割，则与圆周合体而无所失矣。”割得越细，差得越少，割到极限就和圆分毫不差。这已经是非常成熟的极限思想，比欧洲同类想法早了一千多年。

古希腊的阿基米德则认真对待了抛物线山包。他的办法是把山包切成一层又一层越来越小的三角形：第一层一个大三角形，第二层两个小三角形，第三层四个更小的三角形……他证明了每一层三角形的总面积恰好是上一层的四分之一，于是总面积是一个无限相加的和：
\[
S=1+\frac14+\frac1{16}+\frac1{64}+\cdots
\]
他又证明这个无限和等于 $\dfrac43$。也就是说，抛物线山包的面积恰好是内接大三角形的 $\dfrac43$ 倍。在没有代数符号、没有坐标的年代，这几乎是徒手攀岩式的成就。

但阿基米德的方法有一个让人不安的地方：每换一条曲线，就得重新发明一套切割技巧，像每把锁都要单独打一把钥匙。圆的面积有一套办法，抛物线有一套办法，螺线又是一套办法，每一套都依赖只有天才才看得到的巧妙构图。有没有一把\emph{万能钥匙}，对任何曲线都用同一个套路？这个问题卡了数学家将近两千年。

\section{第一次尝试：矩形太少，误差太大}

最自然的尝试是用矩形。把区间 $[-1,1]$ 切成几段，在每一段上立一个矩形，矩形的高取这一段上曲线的高度，把所有矩形面积加起来，当作山包面积的近似。

先切两段试试：左半段上曲线高度在 $0$ 到 $1$ 之间波动，右半段也一样。如果每段都用最高点的高度 $1$，两个矩形面积都是 $1\times 1=1$，总面积 $2$——这就是外接矩形，明显偏大，两个角上多算了一大块。如果每段都用最低点的高度，左段最低点是端点处的 $0$，右段也是 $0$，总面积是 $0$——又小得离谱。一个在 $2$，一个在 $0$，真正的面积藏在中间，可这两个估计之间的差距太大，等于什么都没说。

切成四段会好一些，十段更好，一百段更好。而且这里有一个让人安心的规律：用每段最高点算出的“上和”永远不小于真实面积，用每段最低点算出的“下和”永远不大于真实面积，段数加倍时，上和往下缩、下和往上涨，两者的差距越挤越小。真实面积就像被两堵墙夹在中间的囚犯，墙越砌越近。可只要矩形是有限个，每个矩形的顶边是平的，而曲线的顶是弯的，两堵墙之间就始终留着一条缝。你不能说“取 $100$ 个矩形就是答案”，只能说“取 $100$ 个矩形离答案不远”。那答案本身是什么？难道我们永远只能逼近，永远不能到达？

\begin{fanli}
“矩形越多越准，所以取一百万个矩形，结果就是精确值了。”这个说法看似合理，其实偷换了概念。一百万个矩形仍然是有限个，顶边仍然是平的，得到的仍然是一个近似值，只不过误差很小而已。“近似得非常好”和“精确相等”是两回事。积分思想的关键一步，恰恰是不满足于任何有限的细分，而是让段数趋向无穷，取这些近似值的\emph{极限}——回想第 9 章，极限才是那个“真正到达”的数。
\end{fanli}

\section{把“越切越细”变成一个定义}

走出困境的办法，是承认我们需要的不是一个更大的有限数，而是一个极限过程。把区间 $[a,b]$ 切成 $n$ 小段，每段宽 $\Delta x=\dfrac{b-a}{n}$，在第 $i$ 段里随便取一点，记下曲线在那里的高度 $f(x_i)$，这一段贡献的矩形面积约是 $f(x_i)\,\Delta x$。全部加起来：
\[
S_n=f(x_1)\,\Delta x+f(x_2)\,\Delta x+\cdots+f(x_n)\,\Delta x .
\]
这个和叫作\emph{黎曼和}。当 $n$ 越来越大，$\Delta x$ 越来越小，如果 $S_n$ 趋向一个确定的数 $S$，而且这个数不依赖于你在每段里取哪一点，我们就说函数 $f$ 在 $[a,b]$ 上可积，$S$ 就是它的定积分。

\begin{dingyi}[定积分]
设函数 $f$ 在区间 $[a,b]$ 上有定义。把 $[a,b]$ 等分成 $n$ 段，记 $\Delta x=\dfrac{b-a}{n}$，在每段中任取一点 $x_i$。若当 $n\to\infty$ 时，黎曼和
\[
S_n=\sum_{i=1}^{n} f(x_i)\,\Delta x
\]
的极限存在且与取点方式无关，则称这个极限为 $f$ 在 $[a,b]$ 上的\emph{定积分}，记作
\[
\int_a^b f(x)\,dx=\lim_{n\to\infty}\sum_{i=1}^{n} f(x_i)\,\Delta x .
\]
\end{dingyi}

\begin{zhu}
并不是随手写一个函数都可积。好在本书遇到的函数——连续函数，以及只有少数几个跳跃点的函数——都是可积的。真正把“哪些函数可积”问到底，是大学实变函数课的故事。“与取点方式无关”这个要求也不是摆设：如果换一套取点方式极限就变了，那说明函数太“野”，面积根本没有确定值。
\end{zhu}

\begin{center}
\begin{tikzpicture}[scale=1.1]
  \draw[->] (-0.5,0) -- (5,0) node[right] {$x$};
  \draw[->] (0,-0.4) -- (0,4.6) node[above] {$y$};
  \draw[fill=blue!12] (0,0) rectangle (1,3);
  \draw[fill=blue!12] (1,0) rectangle (2,3.75);
  \draw[fill=blue!12] (2,0) rectangle (3,4);
  \draw[fill=blue!12] (3,0) rectangle (4,3.75);
  \draw[thick,domain=0:4,samples=60] plot (\x,{4-0.25*(\x-2)*(\x-2)});
  \node at (2,-0.35) {\small 用矩形逼近曲边图形的面积};
\end{tikzpicture}
\end{center}

新符号 $\int$ 解决了什么麻烦？它把“切成无穷多份、每份无穷小、再加起来”这一整套极限过程压缩成一个记号。这个拉长的字母 S（sum，求和）提醒我们：积分骨子里就是加法，只不过是加到了极限。$f(x)\,dx$ 则像一个“无穷小矩形”的面积：高是 $f(x)$，宽是 $dx$。

\begin{shuyu}{定积分}
\textbf{直观解释：}曲边图形面积的精确值，由越来越细的矩形面积之和逼近而成。

\textbf{正式定义：}黎曼和在分割无限加细时的极限，记作 $\int_a^b f(x)\,dx$。

\textbf{一个例子：}$\int_0^2 x\,dx=2$，恰好是底为 $2$、高为 $2$ 的三角形面积。
\end{shuyu}

回到里程表的问题。把旅程切成 $n$ 个很短的时间段，第 $i$ 段里的速度近似为 $v(t_i)$，路程约是 $v(t_i)\,\Delta t$，总路程约是黎曼和。令 $n\to\infty$，就得到
\[
\text{总路程}=\int_a^b v(t)\,dt .
\]
在速度—时间图像上，$v(t_i)\,\Delta t$ 是一个细矩形的面积，所以\emph{路程就是速度曲线下方的面积}。速度和路程，一个是“率”，一个是“量”，积分把它们连了起来。如果车是匀速的，速度曲线是水平线，曲线下面积是矩形——这就是小学里“路程等于速度乘时间”；积分做的事情，是把这句小学口诀推广到速度任意变化的情形。

\section{万能钥匙：微积分基本定理}

定义有了，可如果每算一个积分都要真的去求黎曼和的极限，那跟阿基米德每次单独打钥匙也强不了多少。幸运的是，微积分送给我们一份厚礼：计算定积分根本不用求极限，只要找到一个“导数等于 $f$”的函数就够了。

\begin{dingli}[微积分基本定理]
设 $f$ 是 $[a,b]$ 上连续的函数，$F$ 是 $f$ 的一个原函数，即 $F'(x)=f(x)$。那么
\[
\int_a^b f(x)\,dx=F(b)-F(a).
\]
\end{dingli}

这个定理为什么是对的？用路程和速度的语言最清楚。设 $s(t)$ 是里程表读数，那么速度 $v(t)=s'(t)$——这是导数的定义。另一方面，从时刻 $a$ 到时刻 $b$，里程表增加的总路程就是 $s(b)-s(a)$。而我们刚刚用积分算过，总路程等于 $\int_a^b v(t)\,dt$。两个算法算的是同一段路程，所以
\[
\int_a^b v(t)\,dt=s(b)-s(a).
\]
把 $v$ 换成 $f$、把 $s$ 换成 $F$，这就是定理本身。下面把这个直觉写得更像证明一点。

\begin{proof}
把 $[a,b]$ 等分成 $n$ 段，分点为 $a=t_0<t_1<\cdots<t_n=b$，每段长 $\Delta t$。总增量可以拆成一串小增量之和：
\[
F(b)-F(a)=\bigl[F(t_1)-F(t_0)\bigr]+\bigl[F(t_2)-F(t_1)\bigr]+\cdots+\bigl[F(t_n)-F(t_{n-1})\bigr],
\]
右边中间项两两抵消，只剩首尾。每一小段上，由导数的意义，当 $\Delta t$ 很小时
\[
F(t_i)-F(t_{i-1})\approx F'(t_{i-1})\,\Delta t=f(t_{i-1})\,\Delta t ,
\]
所以 $F(b)-F(a)$ 约等于黎曼和 $\sum f(t_{i-1})\,\Delta t$。令 $n\to\infty$，近似变成精确的极限：左边 $F(b)-F(a)$ 不随分割改变，右边趋向 $\int_a^b f(t)\,dt$，于是二者相等。
\end{proof}

注意证明里最关键的一步：一长串括号中间全部抵消，只剩下 $F(b)-F(a)$。积分的“总和”性质和导数的“局部变化率”性质，通过这串抵消严丝合缝地扣在一起。这就是为什么说\emph{微分和积分互为逆运算}。

\begin{zhu}
原函数并不唯一：$x^{2}$ 的导数是 $2x$，$x^{2}+7$ 的导数也是 $2x$，常数项求导就消失了。但这不影响定理——两个原函数只差一个常数，做减法 $F(b)-F(a)$ 时常数自己消掉。所以算定积分时，随便挑一个最顺手的原函数就行。
\end{zhu}

配上数值例子尝尝味道。要算抛物线山包的面积 $\int_{-1}^{1}(1-x^{2})\,dx$，先找原函数：哪个函数的导数是 $1-x^{2}$？试试 $F(x)=x-\dfrac{x^{3}}{3}$，求导验证：$F'(x)=1-x^{2}$，对了。于是
\[
\int_{-1}^{1}(1-x^{2})\,dx=F(1)-F(-1)=\left(1-\frac13\right)-\left(-1+\frac13\right)=\frac43 .
\]
$\dfrac43$！正是阿基米德用无限三角形拼出来的那个数。两千年后，万能钥匙和老锁匠的答案对上了——这是数学里最让人安心的时刻之一。

\begin{toolbox}
本章新增的工具：
\begin{itemize}
  \item \textbf{黎曼和}：把区间切细，用矩形面积之和逼近曲边面积；
  \item \textbf{定积分} $\int_a^b f(x)\,dx$：黎曼和的极限，几何上是曲线下方的面积；
  \item \textbf{微积分基本定理}：找到原函数 $F$，积分等于 $F(b)-F(a)$，导数和积分互为逆运算；
  \item \textbf{切片思想}：立体的体积是截面积函数的积分，截面积处处相等的立体体积相等；
  \item \textbf{级数判别的直觉}：项趋于 $0$ 只是必要条件，不是充分条件；等比级数在公比绝对值小于 $1$ 时收敛。
\end{itemize}
\end{toolbox}

\begin{lianjie}
这条定理把本书的三章缝在了一起：第 9 章的极限给出了积分定义的骨架，第 10 章的导数提供了计算积分的钥匙，本章的积分则把“变化率”还原成“总量”。到第 12 章你会看到，正是对这条定理中“连续”“极限存在”等条件的反复推敲，催生了一整套严格的分析学。
\end{lianjie}

\section{切面包片：从面积到体积}

面积不是积分唯一的用武之地。一只圆锥、一个球，体积怎么算？初中课本给过公式：圆锥体积是同底等高圆柱的三分之一，球体积是 $\dfrac{4}{3}\pi r^{3}$。当时只能背下来，现在可以问一句为什么。

办法还是老一套：切。把球想象成一条面包，垂直于一条水平轴切成薄片。在距球心 $x$ 处下刀（球心记为 $0$，两端是 $-r$ 和 $r$），截面是一个圆。由勾股定理，这个圆的半径平方是 $r^{2}-x^{2}$，面积是 $\pi\left(r^{2}-x^{2}\right)$。每片厚度 $\Delta x$，体积约是 $\pi\left(r^{2}-x^{2}\right)\Delta x$。全部加起来、取极限：
\[
V=\int_{-r}^{r}\pi\left(r^{2}-x^{2}\right)dx .
\]
被积函数只是个二次多项式，找原函数轻而易举：$F(x)=\pi\left(r^{2}x-\dfrac{x^{3}}{3}\right)$。由微积分基本定理，
\[
V=F(r)-F(-r)=\pi\left(r^{3}-\frac{r^{3}}{3}\right)-\pi\left(-r^{3}+\frac{r^{3}}{3}\right)=\frac{4}{3}\pi r^{3} .
\]
课本里那句“背下来就行”，原来是一行积分的结果。

传说阿基米德最得意的发现正是球与圆柱的关系：球的体积恰好是外接圆柱（底半径 $r$、高 $2r$）的三分之二，他要求把这个图形刻在墓碑上。他没有积分符号，却用了“切片比较”的思想——在同样的高度上，球的截面加上圆锥的截面，面积恒等于圆柱的截面——先一步算出了这个比例。无独有偶，中国古代的祖暅（祖冲之之子）把同一思想总结成一句话：“幂势既同，则积不容异。”意思是：两个立体如果在每个等高处的截面积都相等，体积就必然相等。这一原理后来在欧洲以卡瓦列里的名字流传，比祖暅晚了一千多年。切片思想在两个文明里各自生长，说明它确实是逼近体积的“自然规律”。

圆锥体积的“三分之一”也可以这样看：高 $h$、底半径 $r$ 的圆锥，尖朝上放置，在距顶点 $x$ 处的截面半径由相似三角形给出，是 $\dfrac{rx}{h}$，截面面积是 $\pi\dfrac{r^{2}x^{2}}{h^{2}}$。于是
\[
V=\int_0^{h}\pi\frac{r^{2}x^{2}}{h^{2}}\,dx=\frac{\pi r^{2}}{h^{2}}\cdot\frac{h^{3}}{3}=\frac13\pi r^{2}h .
\]
常数 $\dfrac{\pi r^{2}}{h^{2}}$ 提到积分号外面，剩下的 $\int_0^{h}x^{2}dx=\dfrac{h^{3}}{3}$ 正是我们熟悉的那个积分。体积公式不再是记忆负担，而是“体积等于截面积函数的积分”这一句话的推论。

\section{换一个视角：密度、概率和一碗汤}

积分从来不只是“算面积”“算体积”的工具。凡是“一个率累积成一个量”的地方，都有它的影子。看三个例子。

\textbf{质量与密度。}一根不均匀的金属棒，左端稀、右端密，位置 $x$ 处的线密度是 $\rho(x)$（单位长度的质量）。把棒切成 $n$ 小段，第 $i$ 段近似均匀，质量约 $\rho(x_i)\,\Delta x$，加起来取极限：
\[
\text{总质量}=\int_a^b \rho(x)\,dx .
\]
一锅汤里盐的浓度不均匀，整锅汤的含盐量也是浓度对体积的积分——只不过那是三维的积分，大学里会学到。

\textbf{电与功。}电流是电量的变化率，所以一段时间内流过的总电量是电流对时间的积分。举个例子：若电流随时间增大，$I(t)=2t$ 安培，那么前三秒流过的电量是
\[
\int_0^3 2t\,dt=3^{2}-0^{2}=9\ \text{库仑} .
\]
同样，力在一段距离上做的功，等于力对位移的积分。弹簧拉得越长越费劲，力随伸长量变化，用积分才能算清总共做了多少功。

\textbf{概率与密度。}你在车站等车，车在一个小时内的任意时刻到达的可能性相同。车到时刻落在 $10$ 点到 $10$ 点 $15$ 分之间的概率是多少？直觉会答 $\dfrac14$：时间段占整个小时的四分之一。但如果车的到达时间在傍晚更密集呢？这时“可能性”在不同时刻不一样大，需要一条\emph{密度曲线} $p(x)$ 来描述：时刻 $x$ 附近单位时间里的“概率浓度”。那么
\[
P(\text{车在 }c\text{ 到 }d\text{ 之间到达})=\int_c^d p(x)\,dx ,
\]
概率是密度曲线下方的面积，而且全部概率加起来是 $1$：密度曲线在整个范围上的积分等于 $1$。等车这个例子里密度是常数 $\dfrac{1}{60}$（按分钟计），密度曲线是一条水平线，概率就是矩形面积——又回到最朴素的图形。身高、测量误差这类随机量的密度曲线则是中间高、两边低的“钟形”，第 18 章会专门研究它。

\begin{shuyu}{概率密度}
\textbf{直观解释：}一条曲线，它下方某区间上的面积表示随机量落在这个区间里的概率。

\textbf{正式定义：}非负函数 $p(x)$，满足整个范围上的积分等于 $1$；$P(c\le X\le d)=\int_c^d p(x)\,dx$。

\textbf{一个例子：}$p(x)=\dfrac1{60}$（$0\le x\le 60$ 分钟）表示一小时内均匀到达，落在一刻钟内的概率是 $\dfrac{15}{60}=\dfrac14$。
\end{shuyu}

注意概率密度有一个反直觉之处：密度曲线在某一点的高度不是概率。车恰好在 $10$ 点 $30$ 分 $00$ 秒到达的概率是 $0$——一个点的“区间宽度”为零，面积为零。概率只属于区间，密度只是“浓度的快慢”。这和“某时刻的速度不是路程”是完全相同的道理：率是率，量是量，积分才把率攒成量。

\section{无限相加，什么时候有有限的结果}

阿基米德的面积计算里藏着另一个宝藏：$1+\dfrac14+\dfrac1{16}+\cdots=\dfrac43$。无限多个正数相加，结果居然是有限的。这和直觉打架——“无限多个东西加起来不应该无限大吗？”

不一定。第 9 章的芝诺悖论早就埋伏了答案：要跑完 $1$ 公里，先跑 $\dfrac12$ 公里，再跑剩下的 $\dfrac14$，再跑 $\dfrac18$……虽然要经历无限多个阶段，总路程却是有限的：
\[
\frac12+\frac14+\frac18+\cdots=1.
\]
几何上看得更清楚：把一个单位正方形先涂一半，再涂剩下的一半，再涂剩下的一半……每涂一次都更靠近涂满，却永远留一点没涂。涂的总面积趋向 $1$，不会超过。

\begin{center}
\begin{tikzpicture}[scale=1.4]
  \draw (0,0) rectangle (4,4);
  \fill[blue!25] (0,0) rectangle (2,4);
  \fill[blue!40] (2,2) rectangle (4,4);
  \fill[blue!55] (2,1) rectangle (4,2);
  \fill[blue!70] (2,0.5) rectangle (4,1);
  \node at (1,2) {$\frac12$};
  \node at (3,3) {$\frac14$};
  \node at (3,1.5) {$\frac18$};
  \node at (3,0.75) {$\frac1{16}$};
  \node at (2,-0.3) {\small 一半又一半：总面积趋向 $1$};
\end{tikzpicture}
\end{center}

这类每次乘同一个比的和叫\emph{等比级数}。只要公比 $q$ 满足 $|q|<1$，它就有有限的和：
\[
a+aq+aq^{2}+\cdots=\frac{a}{1-q}.
\]
验证一下芝诺的例子：$a=\dfrac12$，$q=\dfrac12$，和为 $\dfrac{1/2}{1-1/2}=1$，分毫不差。这个公式不难证：设 $S=a+aq+\cdots+aq^{n}$，两边乘 $q$ 得 $qS=aq+\cdots+aq^{n+1}$，两式相减，中间项全部抵消，$(1-q)S=a-aq^{n+1}$。当 $|q|<1$ 且 $n\to\infty$ 时 $q^{n+1}\to 0$，于是 $S\to\dfrac{a}{1-q}$。又一次，\emph{抵消}是主角——和微积分基本定理证明里那串抵消遥相呼应。

循环小数里也藏着等比级数。$0.999\ldots$ 到底等不等于 $1$？把它拆开看：
\[
0.999\ldots=0.9+0.09+0.009+\cdots=\frac{0.9}{1-0.1}=1 .
\]
它不是“无限接近 $1$”，而是\emph{就是} $1$，只是换了一件无限长的外衣。极限不是一个永远走不到终点的过程——极限本身就是终点。

有的级数收敛得飞快。最著名的例子藏在自然对数的底 $e$ 里：
\[
e=1+1+\frac12+\frac16+\frac1{24}+\frac1{120}+\cdots
\]
分母是连乘积：$2=1\times2$，$6=1\times2\times3$，$24=1\times2\times3\times4$……每一项都把自己的分母再乘一个大数，所以项缩小的速度比任何等比级数都猛。只加前七项，就能得到 $e\approx2.718$，精确到小数点后两位。银行复利、人口增长、放射性衰变背后都是这个数，而它的庐山真面目，竟是一个收敛极快的无限和。下一章讲函数的幂级数展开时，我们还会回到这个式子。

但是，千万别以为“只要每一项越来越小，总和就有限”。最著名的反例是\emph{调和级数}：
\[
1+\frac12+\frac13+\frac14+\frac15+\cdots
\]
每一项都趋向 $0$，可它的和却是无穷大。

\begin{dingli}[调和级数发散]
调和级数 $1+\dfrac12+\dfrac13+\dfrac14+\cdots$ 的和超过任何给定的数。
\end{dingli}

\begin{proof}
把项分组：第一项 $1$；第二组 $\dfrac12$；第三组 $\dfrac13+\dfrac14$，两项都不小于 $\dfrac14$，所以这一组大于 $\dfrac14+\dfrac14=\dfrac12$；第四组 $\dfrac15+\dfrac16+\dfrac17+\dfrac18$，四项都不小于 $\dfrac18$，所以这一组大于 $4\times\dfrac18=\dfrac12$；依此类推，接下来 $8$ 项都不小于 $\dfrac1{16}$，合起来又大于 $\dfrac12$。每往后翻一倍个数，就能再凑出至少 $\dfrac12$。$\dfrac12$ 加了无限多次，总和当然超过任何界限。
\end{proof}

这个证明最早由中世纪的学者给出，简单得让人意外，结论却深刻：项趋于 $0$ 只是收敛的\emph{必要条件}，远不是充分条件。级数收不收敛，取决于项变小得“够不够快”——所有 $\dfrac1{n^{2}}$ 加起来收敛，而且它等于 $\dfrac{\pi^{2}}{6}$：圆周率竟然从一堆平方数里冒了出来，这个结果是欧拉年轻时算出的，震惊了当时的数学界；而所有 $\dfrac1n$ 加起来发散。那么 $\dfrac1{n^{1.5}}$ 呢？$\dfrac1{n^{1.001}}$ 呢？快与慢之间那条分界线究竟画在哪里，正是分析学最细腻的问题之一。

\begin{fanli}
“每一项都趋于 $0$，所以总和有限。”调和级数当面反驳了这句话。反过来也要小心：总和有限确实能推出项趋于 $0$，但项趋于 $0$ 推不出总和有限。必要条件和充分条件，一字之差，谬以千里。
\end{fanli}

\section{动手逼近一次积分}

\begin{shiyan}
拿一张方格纸，亲手验证微积分基本定理。

\begin{enumerate}
  \item 画出 $y=x^{2}$ 在 $0\le x\le 2$ 的图像（取几个点描出来即可），把 $[0,2]$ 切成 $4$ 段，每段宽 $0.5$，用每段\emph{右端点}的高度立矩形，算面积之和：
  \[
  0.5\times(0.5^{2}+1^{2}+1.5^{2}+2^{2})=0.5\times 7.5=3.75 .
  \]
  再用\emph{左端点}的高度算一次：$0.5\times(0+0.25+1+2.25)=1.75$。真正的面积被夹在这两个数之间。
  \item 用微积分基本定理算精确值：$x^{2}$ 的原函数是 $\dfrac{x^{3}}{3}$，所以
  \[
  \int_0^2 x^{2}\,dx=\frac{2^{3}}{3}-0=\frac83\approx 2.67 .
  \]
  确实落在 $1.75$ 和 $3.75$ 之间。
  \item 把段数加到 $8$ 段（每段 $0.25$），重算左右两个矩形和，把结果填进表里，观察它们怎样从两边挤向 $\dfrac83$：
\end{enumerate}
\begin{center}
\begin{tabular}{cccc}
\toprule
段数 $n$ & 左端点和 & 右端点和 & 精确值 $\frac83$ \\
\midrule
$4$ & $1.75$ & $3.75$ & $2.666\ldots$ \\
$8$ & $2.19$ & $3.19$ & $2.666\ldots$ \\
\bottomrule
\end{tabular}
\end{center}
有余力的话用计算器算 $n=16$，看看两个近似值各前进了多少，猜测误差大约按什么速度缩小。
\end{shiyan}

\begin{pandeng}
想继续往前走的读者，可以留意这些关键词：\emph{换元积分法}与\emph{分部积分法}（计算积分的两把扳手）；\emph{反常积分}（区间无限长或被积函数爆炸时怎么办，例如 $\int_1^{\infty}\frac{1}{x^{2}}dx=1$）；\emph{数值积分}（计算机如何用梯形法、抛物线法高效逼近积分）；\emph{黎曼 $\zeta$ 函数}（级数 $\sum\frac{1}{n^{s}}$ 收敛与发散的分界线通向第 15 章的素数故事）；以及大学里更一般的\emph{勒贝格积分}——它把“竖着切条”改成“横着切层”，能处理更怪的函数。
\end{pandeng}

\section*{思考题}

\textbf{观察题}
\begin{sikao}
  \item 一辆车的速度表读数始终是正的，但司机先向前开、又倒车回到起点。速度曲线下的“面积”和“总路程”还是一回事吗？提示：倒车时速度该记成正还是负？积分值和里程表实际磨损的路各对应什么？
  \item 观察水表、电表、手机流量统计：它们各自是“率”还是“量”？和它们对应的另一个量是什么？提示：每一对“率与量”之间都藏着一个积分和一个导数。
\end{sikao}

\textbf{证明题}
\begin{sikao}
  \item 不查公式，证明 $\dfrac13+\dfrac19+\dfrac1{27}+\cdots=\dfrac12$。提示：模仿正文中等比级数的“乘 $q$ 再相减”的办法，或者画一个正方形反复取三分之一。
  \item 用微积分基本定理解释：为什么 $\int_a^b x\,dx=\dfrac{b^{2}-a^{2}}{2}$ 恰好等于一个梯形的面积？提示：$y=x$ 的图像在 $[a,b]$ 上围出的是什么图形？梯形的两条底边各是多长？
\end{sikao}

\textbf{挑战题}
\begin{sikao}
  \item 调和级数发散，但把其中所有分母含数字 $9$ 的项（如 $\dfrac19$、$\dfrac1{19}$、$\dfrac1{92}$）统统删掉后，剩下的和竟然收敛。你能估计它为什么收敛吗？提示：一位数、两位数、三位数……各有多少个不含 $9$ 的数？它们各自的倒数和大约被什么控制住？把各段的界加成等比级数试试。
  \item 正文证明了 $\int_a^b f(x)\,dx=F(b)-F(a)$ 依赖“$F$ 的导数是 $f$”。如果只知道 $f$ 连续，能不能\emph{造出}一个原函数？提示：考虑函数 $F(x)=\int_a^x f(t)\,dt$，估计 $\dfrac{F(x+h)-F(x)}{h}$ 在 $h$ 很小时是多少——这正是下一章要严格化的思想。
\end{sikao}

\section*{通往下一章的门}

这一章我们走得很快，快到鞋底有点打滑。黎曼和的极限“存在”是什么意思？凭什么说上和与下和两堵墙越砌越近，最后就一定严丝合缝地挤出一个数，而不是永远留着一条缝？微积分基本定理要求 $f$“连续”，可什么叫连续——“笔不离纸”的画法能当成定义吗？调和级数发散、$\sum\frac{1}{n^{2}}$ 收敛，分界线背后有没有一把统一的尺子？

这些问题都不难懂，却难得漂亮。为了回答它们，十九世纪的数学家们把“无限接近”“连续”“收敛”这些日常词汇逐一锻造成了精确的定义，建起了一门叫\emph{分析}的学问。下一章我们就走进这座锻造车间，看看直觉是怎样被炼成定理的——并且在那里重逢一个老朋友：用无穷多个简单的波，拼出任何复杂的声音。
